% Part: lambda-calculus
% Chapter: syntax
% Section: alpha

\documentclass[../../../include/open-logic-section]{subfiles}

\begin{document}

\olfileid[hi]{lam}{syn}{alp}
\olsection{$\alpha$-परिवर्तन}

$\lambd[x][x]$ और $\lambd[y][y]$ के बीच क्या संबंध है? दोनों तादात्म्य
फलन को निरूपित करते हैं। निश्चय ही, वाक्यविन्यास की दृष्टि से वे भिन्न पद
हैं। वे केवल आबद्ध चर के नाम में भिन्न हैं, और एक दूसरे के आबद्ध चर का
``पुनर्नामकरण'' करने से मिलता है। इसे \emph{$\alpha$-परिवर्तन} कहते हैं।

\begin{defn}[आबद्ध चर का परिवर्तन, $\aconvone$]
  यदि किसी पद $M$ में $\lambd[x][N]$ का कोई आविर्भाव है, $y \notin
  \FV{N}$, और $\Subst{N}{y}{x}$ परिभाषित है, तो इस आविर्भाव को
  \begin{equation*}
    \lambd[y][\Subst{N}{y}{x}]
  \end{equation*}
  से बदलकर $M'$ प्राप्त करने को \emph{आबद्ध चर का परिवर्तन} कहते हैं, और
  इसे $M \redone[\alpha] M'$ लिखते हैं।
\end{defn}

\begin{defn}[संबंध की संगतता]
  पदों पर कोई संबंध $R$ \emph{संगत} कहलाता है यदि वह निम्न शर्तें संतुष्ट
  करे:
  \begin{enumerate}
  \item यदि $R N N'$, तो $R \lambd[x][N] \lambd[x][N']$।
  \item यदि $R P P'$, तो $R (PQ) (P'Q)$।
  \item यदि $R Q Q'$, तो $R (PQ) (PQ')$।
  \end{enumerate}
\end{defn}

अतः परिभाषा को फिर से कहें:
\begin{defn}[आबद्ध चर का परिवर्तन, $\aconvone$]
  \emph{आबद्ध चर का परिवर्तन} ($\redone[\alpha]$) पदों पर वह सबसे छोटा
  संगत संबंध है जो निम्न शर्त संतुष्ट करता है:
  \begin{align*}
    &\lambd[x][N] \redone[\alpha] \lambd[y][\Subst{N}{y}{x}] && \text{यदि
      $x \neq y$, $y \notin \FV{N}$} \\
    & &&\text{और $\Subst{N}{y}{x}$ परिभाषित है।}
  \end{align*}
\end{defn}

यहाँ ``सबसे छोटा'' का अर्थ है कि संबंध में केवल वही युग्म हैं जिनकी संगतता
और अतिरिक्त शर्त से आवश्यकता पड़ती है, और कोई नहीं। अतः इस संबंध को इस
प्रकार भी परिभाषित किया जा सकता है:

\begin{defn}[आबद्ध चर का परिवर्तन, $\aconvone$] \ollabel{defn:aconvone}
  \emph{आबद्ध चर का परिवर्तन} ($\aconvone$) आगमनतः इस प्रकार परिभाषित है:
  \begin{enumerate}
  \item यदि $N \aconvone N'$, तो $\lambd[x][N] \aconvone
    \lambd[x][N']$। \ollabel{defn:aconvone1}
  \item यदि $P \aconvone P'$, तो $(PQ) \aconvone (P'Q)$। \ollabel{defn:aconvone2}
  \item यदि $Q \aconvone Q'$, तो $(PQ) \aconvone (PQ')$। \ollabel{defn:aconvone3}
  \item यदि $x \neq y$, $y \notin \FV{N}$ और $\Subst{N}{y}{x}$ परिभाषित है, तो
    $\lambd[x][N] \redone[\alpha] \lambd[y][\Subst{N}{y}{x}]$.
    \ollabel{defn:aconvone4}
  \end{enumerate}
\end{defn}

ये परिभाषाएँ तुल्य हैं, किंतु प्रमाण अभ्यास के लिए छोड़ते हैं। अब से हम
आगमनात्मक परिभाषा उपयोग करेंगे।

\begin{defn}[$\alpha$-परिवर्तन, $\aconv$]
  \emph{$\alpha$-परिवर्तन} ($\aconv$) पदों पर वह सबसे छोटा परावर्ती और
  संक्रामक संबंध है जिसमें $\aconvone$ समाहित है।
\end{defn}

ऊपर की तरह, ``सबसे छोटा'' का अर्थ है कि संबंध में केवल संक्रामकता और
$\aconvone$ से आवश्यक युग्म हैं। इससे निम्न तुल्य परिभाषा मिलती है:

\begin{defn}[$\alpha$-परिवर्तन, $\aconv$] \ollabel{defn:aconv}
  \emph{$\alpha$-परिवर्तन} ($\aconv$) आगमनतः इस प्रकार परिभाषित है:
  \begin{enumerate}
  \item यदि $P \aconv Q$ और $Q \aconv R$, तो $P \aconv R$।
    \ollabel{defn:aconv1}
  \item यदि $P \aconvone Q$, तो $P \aconv Q$। \ollabel{defn:aconv2}
  \item $P \aconv P$. \ollabel{defn:aconv3}
  \end{enumerate}
\end{defn}

\begin{ex}
  $\lambd[x][f x]$, ~$\alpha$-परिवर्तित होकर $\lambd[y][f y]$ बनता है,
  और उलटा भी। अनौपचारिक रूप से, दोनों ऐसे फलन हैं जो कोई तर्क स्वीकार करके
  उस तर्क पर $f$ का मान लौटाते हैं और परिवेश-चर~$f$ को संदर्भित करते हैं।

  $\lambd[x][f x]$, ~$\alpha$-परिवर्तित होकर $\lambd[x][g x]$ नहीं बनता।
  अनौपचारिक रूप से, वे क्रमशः परिवेश-चरों $f$ और~$g$ को संदर्भित करते हैं,
  इसलिए वे भिन्न फलन हैं: जिन परिवेशों में $f$ और $g$ भिन्न हों उनमें उनका
  व्यवहार भी भिन्न होता है।
\end{ex}

\begin{prob}
  क्या पदों के निम्न युग्म $\alpha$-परिवर्तनीय हैं?
  \begin{enumerate}
  \item $\lambd[x][\lambd[y][x]]$ और $\lambd[y][\lambd[x][y]]$
  \item $\lambd[x][\lambd[y][x]]$ और $\lambd[c][\lambd[b][a]]$
  \item $\lambd[x][\lambd[y][x]]$ और $\lambd[c][\lambd[b][a]]$
  \end{enumerate}
\end{prob}

\begin{lem}\ollabel{lem:fv-one}
  यदि $P \aconvone Q$, तो $\FV{P} = \FV{Q}$।
\end{lem}

\begin{proof}
  $P \aconvone Q$ के !!{derivation} पर आगमन द्वारा।
  \begin{enumerate}
  \item यदि अंतिम नियम \olref{defn:aconvone4} है, तो $P$,
    ~$\lambd[x][N]$ रूप का और $Q$, ~$\lambd[y][\Subst{N}{y}{x}]$ रूप का
    है, जहाँ $x \neq y$, $y \notin \FV{N}$ और $\Subst{N}{y}{x}$ परिभाषित
    है। हम इस आधार पर स्थितियाँ अलग करते हैं कि $x \in \FV{N}$ है या नहीं:
    \begin{enumerate}
    \item यदि $x \in FV(N)$, तो:
      \begin{align*}
        \FV{\lambd[y][\Subst{N}{y}{x}]} &=
        \FV{\Subst{N}{y}{x}} \setminus \{y\} \\
        & = ((\FV{N} \setminus \{x\}) \cup \{y\}) \setminus \{y\}
         && \text{\olref[sub]{thm:infv} से} \\
        & = \FV{N} \setminus \{x\} \\
        & = \FV{\lambd[x][N]}
      \end{align*}
    \item यदि $x \notin FV(N)$, तो:
      \begin{align*}
        \FV{\lambd[y][\Subst{N}{y}{x}]}
        & = FV{\Subst{N}{y}{x}} \setminus \{y\} \\
        & = \FV{N} \setminus \{x\}
         && \text{\olref[sub]{thm:notinfv} से} \\
        & = \FV{\lambd[x][N]}.
      \end{align*}
    \end{enumerate}
  \item अन्य तीन स्थितियाँ अभ्यास के लिए छोड़ी गई हैं।
  \end{enumerate}
\end{proof}

\begin{prob}
  \olref[lam][syn][alp]{lem:fv-one} का प्रमाण पूरा करें।
\end{prob}

\begin{lem}\ollabel{lem:inv}
  यदि $P \aconvone Q$, तो $Q \aconvone P$।
\end{lem}

\begin{proof}
  $P \aconvone Q$ के !!{derivation} पर आगमन।
  \begin{enumerate}
  \item यदि अंतिम नियम \olref{defn:aconvone4} है, तो $P$,
    ~$\lambd[x][N]$ रूप का और $Q$, ~$\lambd[y][\Subst{N}{y}{x}]$ रूप का
    है, जहाँ $x \neq y$, $y \notin \FV{N}$ और $\Subst{N}{y}{x}$
    परिभाषित है। पहले, \olref[sub]{thm:clr} से $y \notin
    \FV{\Subst{N}{y}{x}}$। \olref[sub]{thm:inv} से
    $\Subst{\Subst{N}{y}{x}}{x}{y}$ न केवल परिभाषित है, बल्कि~$N$ के बराबर
    भी है। तब \olref{defn:aconvone4} से
    $\lambd[y][\Subst{N}{y}{x}] \aconvone
    \lambd[x][\Subst{\Subst{N}{y}{x}}{x}{y}] = \lambd[x][N]$।
  \end{enumerate}
\end{proof}

\begin{prob}
  \olref[lam][syn][alp]{lem:inv} का प्रमाण पूरा करें।
\end{prob}

\begin{thm}
  $\alpha$-परिवर्तन पदों पर तुल्यता-संबंध है, अर्थात् वह परावर्ती, सममित
  और संक्रामक है।
\end{thm}

\begin{proof}
  \begin{enumerate}
  \item प्रत्येक पद $M$ के लिए आबद्ध चरों में \emph{शून्य} परिवर्तन करके
    $M$ को $M$ में बदला जा सकता है।
  \item यदि आबद्ध चरों में परिवर्तनों के किसी अनुक्रम द्वारा $P$,
    ~$\alpha$-परिवर्तित होकर $Q$ बनता है, तो \olref{lem:inv} से उन
    परिवर्तनों को उलटे क्रम में प्रतिलोम करके~$Q$ से~$P$ प्राप्त कर सकते
    हैं।
  \item यदि आबद्ध चरों में परिवर्तनों के किसी अनुक्रम द्वारा $P$,
    ~$\alpha$-परिवर्तित होकर $Q$ बनता है, और किसी दूसरे अनुक्रम द्वारा $Q$
    बदलकर $R$ बनता है, तो पहले पहला और फिर दूसरा अनुक्रम लागू करके $P$ को
    $R$ में बदल सकते हैं।
  \end{enumerate}
\end{proof}

अब से हम कहेंगे कि $M$ और $N$ \emph{$\alpha$-तुल्य} हैं, $M \aeq N$, तभी
और केवल तभी जब $M$, ~$\alpha$-परिवर्तित होकर $N$ बनता है (और जैसा अभी
दिखाया, ऐसा तभी और केवल तभी है जब $N$, ~$\alpha$-परिवर्तित होकर~$M$ बनता
है)।

\begin{thm}\ollabel{thm:fv}
  यदि $M \aeq N$, तो $\FV{M} = \FV{N}$।
\end{thm}

\begin{proof}
  \olref{lem:fv-one} से तात्कालिक।
\end{proof}

\begin{lem}\ollabel{lem:sub:R}
  यदि $R \aeq R'$ और $\Subst{M}{R}{y}$ परिभाषित है, तो
  $\Subst{M}{R'}{y}$ परिभाषित है और $\Subst{M}{R}{y}$ के $\alpha$-तुल्य
  है।
\end{lem}

\begin{proof}
  अभ्यास।
\end{proof}

\begin{prob}
  \olref[lam][syn][alp]{lem:sub:R} सिद्ध करें।
\end{prob}

स्मरण करें कि \olref[sub]{sec} में कुछ स्थितियों में प्रतिस्थापन अपरिभाषित
है; फिर भी, पदों पर $\alpha$-परिवर्तन का उपयोग करके हम आबद्ध चरों का
पुनर्नामकरण कर सकते हैं और प्रतिस्थापन को सदा परिभाषित बना सकते हैं। यह
परिणाम $\alpha$-तुल्यता को सुरक्षित रखता है, जैसा अगला प्रमेय दिखाता है:

\begin{thm}\ollabel{thm:sub}
  किन्हीं $M$, $R$ और~$y$ के लिए कोई ऐसा $M'$ है कि $M \aeq M'$ और
  $\Subst{M'}{R}{y}$ परिभाषित है। इसके अतिरिक्त, यदि कोई दूसरा युग्म
  $M'' \aeq M$ और $R''$ ऐसा है कि $\Subst{M''}{R''}{y}$ परिभाषित है और
  $R'' \aeq R$, तो $\Subst{M'}{R}{y} \aeq \Subst{M''}{R''}{y}$।
\end{thm}

\begin{proof}
  ~$M$ के निर्माण पर आगमन द्वारा:
  \begin{enumerate}
  \item $M$ कोई चर~$z$ है: अभ्यास।
  \item मान लें $M$, ~$\lambd[x][N]$ रूप का है। $x$ और $y$ से भिन्न ऐसा
    चर $z$ चुनें कि $z \notin \FV{N}$ और $z \notin \FV{R}$।
    आगमन-परिकल्पना से कोई ऐसा $N'$ है कि $N' \aeq N$ और
    $\Subst{N'}{z}{x}$ परिभाषित है। तब
    \olref{defn:aconvone}\olref{defn:aconvone1} से $\lambd[x][N] \aeq
    \lambd[x][N']$ भी। अब
    \olref{defn:aconvone}\olref{defn:aconvone4} से $\lambd[x][N'] \aeq
    \lambd[z][\Subst{N'}{z}{x}]$। हम ऐसा कर सकते हैं क्योंकि $z \ne x$,
    $z \notin FV(N')$ और $\Subst{N'}{z}{x}$ परिभाषित है। अंततः
    $\Subst{\lambd[z][\Subst{N'}{z}{x}]}{R}{y}$ परिभाषित है, क्योंकि $z
    \neq y$ और $z \notin FV(R)$।

    इसके अतिरिक्त, यदि इन्हीं शर्तों को संतुष्ट करने वाले कोई अन्य $N''$
    और $R''$ हों, तो
    \begin{multline*}
      \Subst{(\lambd[z][\Subst{N''}{z}{x}])}{R''}{y} =\\
    \begin{aligned}
      &= \lambd[z][\Subst{\Subst{N''}{z}{x}}{R''}{y}] \\
      &= \lambd[z][\Subst{\Subst{N''}{z}{x}}{R}{y}] && \text{
        \olref{lem:sub:R} से}\\
      &=\lambd[z][\Subst{\Subst{N'}{z}{x}}{R}{y}]
      && \text{आगमन-परिकल्पना से}\\
      &=\Subst{(\lambd[z][\Subst{N'}{z}{x}])}{R}{y}
    \end{aligned}
    \end{multline*}
    \item $M$, ~$(PQ)$ रूप का है: अभ्यास।
  \end{enumerate}
\end{proof}

\begin{prob}
  \olref[lam][syn][alp]{thm:sub} का प्रमाण पूरा करें।
\end{prob}

\begin{cor}\ollabel{cor:sub}
  किन्हीं $M$, $R$ और~$y$ के लिए $M'$ और $R'$ का कोई ऐसा युग्म है कि
  $M' \aeq M$, $R \aeq R'$ और $\Subst{M'}{R'}{y}$ परिभाषित है। इसके
  अतिरिक्त, यदि कोई दूसरा युग्म $M'' \aeq M$ और $R''$ ऐसा है कि
  $\Subst{M'}{R'}{y}$ परिभाषित है, तो $\Subst{M'}{R'}{y} \aeq
  \Subst{M''}{R''}{y}$।
\end{cor}

\begin{proof}
  \olref{thm:sub} से तात्कालिक।
\end{proof}

\end{document}

